% =====================================================================
\chapter{Formal Verification with ProVerif}\label{ch:pv}
% =====================================================================
The proof sketches of Chapter~\ref{sec:security} argue each property separately. A protocol can still fail through an interaction between steps that no
single proof considers. We therefore modelled the framework, and the base scheme's credential, in ProVerif~2.05~\cite{proverif,proverifbook}. ProVerif
explores all interleavings of an unbounded number of sessions against a Dolev--Yao attacker. All models are in \texttt{final/proverif/}, and the full
output of every run is kept next to them (\texttt{*.out}).

\section{Modelling choices}
\begin{itemize}
  \item \textbf{Certificateless signatures.} The partial key is a function $\mathsf{partial}(s,(\mathit{id},X))$ of the master key and of the identity
        \emph{and} the entity's public value, mirroring $h=H_1(\PID,\ep,X,Y)$. A signature $\mathsf{clsign}(m,x,d)$ needs both the secret value $x$ and the partial
        key $d$. Verification $\mathsf{clcheck}$ needs only $\mathit{id}$, $X=\mathsf{pubx}(x)$ and $\Ppub=\mathsf{pubm}(s)$.
  \item \textbf{Diffie--Hellman} is modelled with the standard equation $\mathsf{exp}(\mathsf{exp}(g,a),b)=\mathsf{exp}(\mathsf{exp}(g,b),a)$; KDF, AEAD and MAC
        are modelled as perfect.
  \item \textbf{Type-I vs Type-II.} In the Type-I model the master key $s$ is private. In the Type-II model $s$ is public, so the attacker \emph{is} the TA.
  \item \textbf{Forward secrecy} uses a second \emph{phase} in which every long-term secret (and, for C1, even $s$) is published. Secrets sent in phase~0
        must stay secret.
  \item \textbf{Single-use tickets} (C3) are modelled with a ProVerif table of used nonces.
  \item \textbf{The base scheme's credential} is abstracted as $\Gamma=\mathsf{gam}(s_i,\xi_i,R_{tu},Q_{tu})$, where every argument except $s_i$ is public, exactly
        as in equation~\eqref{eq:gamma2}.
\end{itemize}

\section{Results}
\begin{table}[H]
\centering\small
\caption{ProVerif results. ``Expected'' is what the design predicts; every result matches it.}
\label{tab:pv}
\begin{tabularx}{\linewidth}{@{}l l >{\raggedright\arraybackslash}X c c@{}}
\toprule
\# & Model & Property & Expected & Result\\
\midrule
Q1 & C1, Type I & TUAV accepts UAV's registered key $\Rightarrow$ UAV sent it (injective) & true & \yes\ true\\
Q2 & C1, Type I & TUAV accepts the UAV's identity with \emph{any} key $\Rightarrow$ UAV sent it & true & \yes\ true\\
Q3 & C1, Type I & UAV obtains a key $\Rightarrow$ TUAV sent it to that UAV (injective) & true & \yes\ true\\
Q4 & C1, Type I & secrecy of the group key $\GK_a$ & true & \yes\ true\\
Q5 & C1, Type I & data under $\GK_a$ stays secret after \emph{all} long-term keys leak & true & \yes\ true\\
\midrule
Q6 & C1, Type II & malicious TA: accept under the \emph{registered} key $\Rightarrow$ UAV sent it & true & \yes\ true\\
Q7 & C1, Type II & malicious TA: accept the identity with \emph{any} key $\Rightarrow$ UAV sent it & false & \halfmark\ false\\
Q8 & C1, Type II & data under the pairwise key $K_i$ stays secret from the TA & true & \yes\ true\\
Q9 & C1, Type II & group key secret from the TA & false & \halfmark\ false\\
\midrule
Q10 & Rekey & revoked, compromised member knows the \emph{old} key (sanity) & known & \yes\ known\\
Q11 & Rekey & revoked member cannot learn the \emph{new} key & true & \yes\ true\\
\midrule
Q12 & C3 & new TUAV accepts $\M$ $\Rightarrow$ $\M$ started that handover (injective) & true & \halfmark\ non-inj.\ only\\
Q13 & C3 & $\M$ completes $\Rightarrow$ new TUAV sent that response (injective) & true & \yes\ true\\
Q14 & C3 & secrecy of the new cell's group key & true & \yes\ true\\
Q15 & C3 & data in the new cell stays secret after $K_{ab}$ and $K_\M$ leak & true & \yes\ true\\
\midrule
Q16 & Base scheme & TUAV/TA accept UAV $\Rightarrow$ UAV sent it (TA's database leaked) & false & \no\ \textbf{attack}\\
\bottomrule
\end{tabularx}
\end{table}

\section{Interpretation}
\paragraph{The base scheme (Q16).} ProVerif finds an attack in which the attacker, knowing $s_i$ from the TA's database, produces a credential the
verifier accepts for a UAV that never sent anything. This is Finding~\ref{f:escrow}, found automatically.

\paragraph{C1 against outsiders (Q1--Q5).} All properties hold, including mutual authentication in both directions and forward secrecy. Q5 holds even
when the master key $s$ is published in phase~1, because session data depends on erased ephemeral exponents.

\paragraph{C1 against a malicious TA (Q6--Q9): the exact boundary.} This is the most informative result.
\begin{itemize}
  \item \textbf{Q6 holds}: the TA cannot sign under a UAV's \emph{registered} key, because it lacks $x$. This is the escrow-freeness the base scheme lacks.
  \item \textbf{Q8 holds}: the TA cannot read data protected by the UAV's pairwise key $K_i$.
  \item \textbf{Q7 fails} (as predicted): the TA can create a \emph{second} key $X'$ for the same pseudonym, issue $\mathsf{partial}(s,(\mathit{id},X'))$, and be
        accepted. This is key replacement, which every certificateless scheme permits (Girault level~2).
  \item \textbf{Q9 fails} (as predicted): the TA can enrol a phantom member and so receive the group key.
\end{itemize}
These four results replace a vague claim (``the TA cannot impersonate'') with a precise one: \emph{the TA cannot impersonate a registered key or read pairwise
traffic; it can create detectable new keys}. Section~\ref{sec:security} and the limitations in Chapter~\ref{ch:conclusion} are worded accordingly.

\paragraph{Revocation (Q10--Q11).} After a member is removed and fully compromised, the new key stays secret. The base scheme fails this in the
computational setting (Finding~\ref{f:crt}, reproduced in code). The polynomial-root attack is algebraic and cannot be expressed in ProVerif's symbolic
model, which is why we demonstrate it with code rather than ProVerif.

\paragraph{Handover (Q12--Q15).} Secrecy and forward secrecy hold, and so does the member-side injective agreement (Q13). For Q12, ProVerif proves
non-injective agreement but cannot prove injectivity. The reason is that in ProVerif a table lookup and the later insert are not atomic, so two parallel
copies of the new TUAV can both accept one replayed ticket before either records the nonce. In our implementation the TUAV processes handover requests
sequentially, so check-and-insert is atomic and the replay is rejected (test \texttt{test\_c3\_handover}). We report the ProVerif result as it is. Even in the
abstract model, a replayed ticket gives the attacker nothing: by Q13--Q15 it cannot complete the session or learn any key.

\begin{limitbox}
ProVerif works in the symbolic model: primitives are perfect and the attacker cannot exploit algebra beyond the stated equations. Its results complement,
and do not replace, the computational arguments of Chapter~\ref{sec:security}. The models abstract timestamps, epochs and the batch equation; the batch
test's soundness rests on Theorem~\ref{thm:batch} and is checked in code (experiment E7).
\end{limitbox}
